\section{Theory}
\label{sec:Theorie}




\subsection{Uncertainties of the measurement}
\label{sec:fehlerrechnung}
%\subsection{Mittelwert}
The mean value is calculated with
\begin{equation}
\bar{x} = \frac{1}{N} \sum_{i=1}^{N} x_i \, .
\label{eqn:mittelwert}
\end{equation}
%\subsection{Standardabweichung}
\\
The standard deviation is given by
\begin{equation}
\sigma = \sqrt{\frac{1}{N - 1} \sum_{i=1}^{N} \left( x_i - \bar{x}_N \right)^2} \, .
\label{eqn:standardabweichung}
\end{equation}
%%\subsubsection{Standardabweichung des Mittelwerts}
The standard deviation of the mean value is given by
\begin{equation}
\Delta \bar{x} = \frac{\sigma}{\sqrt{N}} \, .
\label{eqn:standardmittelwert}
\end{equation}
%%\subsection{Gaußsche Fehlerfortpflanzung}
The deviation of a value derived from experimental parameters can be calculated using the Gaussian error propagation
\begin{equation}
\Delta f = \sqrt{\sum_{i=1}^{N} \left( \frac{\partial f}{\partial x_i} \right)^2 \cdot (\Delta x_i)^2} \, .
\label{eqn:fehlerfortpflanzung}
\end{equation}
The relative deviation of an experimental parameter from the theoretical parameter is given by ~\cite{relabw}
\begin{equation}
    \increment x_{\mathrm{rel}} = \left| \frac{x_{\mathrm{\exp}}-x_{\mathrm{theo}}}{x_{\mathrm{theo}}}\right|\, .
    \label{eqn:relabw}
\end{equation}

